\section{数据洞察与评估仪表板}
\subsection{核心仪表板视图}
Graham 建议构建一套 Agent 评估仪表板，用以实时监控成功率、干预次数、平均运行时长与安全门控触发记录。仪表板的每一项指标都应该关联具体的事件 log，以便在指标异常时追溯根因。

\begin{table}[H]
\centering
\begin{tabular}{p{0.28\textwidth} p{0.42\textwidth} p{0.28\textwidth}}
\toprule
指标 & 说明 & 目标 \\
\midrule
通过率 & Agent 执行 patch 并通过测试的比例 & \(\geq 40\%\)（Lite）、\(\geq 20\%\)（Full） \\
\midrule
人工干预率 & 需要人类暂停或审批的比例 & \(\leq 10\%\) \\
\midrule
运行时长 & 平均从 Plan 到 Verify 完成的时间 & \(\leq 20 \text{ 分钟}\) \\
\midrule
安全触发 & 涉及数据库/部署的命令数量 & 需人工审核 \\
\bottomrule
\end{tabular}
\caption{Agent 仪表板中关注的核心指标。}
\end{table}

\begin{knowledgebox}{仪表板数据依赖}
每个指标都应与事件 log（命令、diff、测试输出）关联，便于回溯并训练可解释的模型。
\end{knowledgebox}

\subsection{仪表板反馈循环}
仪表板不仅监控，还应支持“反馈 -> 调整”流程：当某个指标出现回落，系统自动把失败 log 发送给 Agent 团队，让他们出具新的 prompt/工具组合，并把新策略推送到测试环境。

\begin{importantbox}{数据反馈实践}
设计自动化报告：指标异常（如通过率<30\%）时，自动打开 issue、附上 diff/log 并提醒团队，防止问题长期沉没。
\end{importantbox}

\subsection{本章小结}
\begin{itemize}
    \item 实时仪表板是判断 Agent 成熟度的关键窗户。
    \item 把异常事件变成自动反馈循环，可以加速策略迭代。
\end{itemize}

